\documentclass[11pt, a4paper]{article}

% Gemeinsame Style-Definitionen (TEXINPUTS muss templates/ enthalten — siehe scripts/build_tex.py)
% bfw-style.tex — gemeinsame Style-Definitionen für alle Handouts/Aufgabenhefte
% Voraussetzung: Kompilation mit xelatex (wegen fontspec)

\usepackage[ngerman]{babel}
\usepackage{geometry}
\geometry{a4paper, top=2.2cm, bottom=2.2cm, left=2.3cm, right=2.3cm}

\usepackage{fontspec}
% Fallback-Kette: Source Sans Pro -> Calibri -> Segoe UI -> Latin Modern Sans
% (Latin Modern ist immer mit MiKTeX vorhanden)
\IfFontExistsTF{Source Sans Pro}{%
  \setmainfont{Source Sans Pro}[Ligatures=TeX]%
}{\IfFontExistsTF{Calibri}{%
  \setmainfont{Calibri}[Ligatures=TeX]%
}{\IfFontExistsTF{Segoe UI}{%
  \setmainfont{Segoe UI}[Ligatures=TeX]%
}{%
  \setmainfont{Latin Modern Sans}[Ligatures=TeX]%
}}}
\IfFontExistsTF{Source Code Pro}{%
  \setmonofont{Source Code Pro}[Scale=0.92]%
}{\IfFontExistsTF{Consolas}{%
  \setmonofont{Consolas}[Scale=0.92]%
}{%
  \setmonofont{Latin Modern Mono}[Scale=0.92]%
}}

\usepackage{xcolor}
% Farbpalette: hell, freundlich, modern
\definecolor{bfwPrimary}{HTML}{0EA5A4}    % Petrol/Türkis - Primär
\definecolor{bfwPrimaryDark}{HTML}{0F766E} % dunkler Petrol für Headings
\definecolor{bfwAccent}{HTML}{F59E0B}     % Amber - Akzent
\definecolor{bfwSuccess}{HTML}{10B981}    % Grün - Erfolg / Lösung
\definecolor{bfwWarn}{HTML}{F97316}       % Orange - Achtung
\definecolor{bfwInk}{HTML}{0F172A}        % fast Schwarz
\definecolor{bfwInkSoft}{HTML}{475569}    % gedämpftes Grau
\definecolor{bfwBgSoft}{HTML}{F8FAFC}     % off-white
\definecolor{bfwBgTeal}{HTML}{ECFEFF}     % helles Türkis
\definecolor{bfwBgAmber}{HTML}{FEF3C7}    % helles Gelb
\definecolor{bfwBgRose}{HTML}{FFE4E6}     % helles Rosé für Eselsbrücken

\usepackage[colorlinks=true, linkcolor=bfwPrimaryDark, urlcolor=bfwPrimaryDark]{hyperref}
\usepackage{enumitem}
\setlist[itemize]{leftmargin=*, itemsep=2pt, topsep=2pt}
\setlist[enumerate]{leftmargin=*, itemsep=2pt, topsep=2pt}

\usepackage[most]{tcolorbox}
\usepackage{booktabs}
\usepackage{tikz}
\usetikzlibrary{decorations.pathmorphing, positioning, shapes.geometric, arrows.meta}

% Merkkästchen — wichtige Definition / Kernpunkt
\newtcolorbox{merksatz}[1][]{%
  enhanced, breakable,
  colback=bfwBgTeal,
  colframe=bfwPrimary,
  boxrule=0.6pt,
  arc=4pt,
  left=10pt, right=10pt, top=8pt, bottom=8pt,
  fonttitle=\bfseries\color{bfwPrimaryDark},
  title={\faLightbulb\ Merksatz},
  #1
}

% Eselsbrücke — Mnemoniks
\newtcolorbox{eselsbruecke}[1][]{%
  enhanced, breakable,
  colback=bfwBgRose,
  colframe=bfwAccent,
  boxrule=0.6pt,
  arc=4pt,
  left=10pt, right=10pt, top=8pt, bottom=8pt,
  fonttitle=\bfseries\color{bfwWarn},
  title={Eselsbrücke},
  #1
}

% Beispiel-Box
\newtcolorbox{beispiel}[1][]{%
  enhanced, breakable,
  colback=bfwBgSoft,
  colframe=bfwInkSoft,
  boxrule=0.4pt,
  arc=3pt,
  left=10pt, right=10pt, top=8pt, bottom=8pt,
  fonttitle=\bfseries\color{bfwInk},
  title={Beispiel},
  #1
}

% Lösungs-Box (für Aufgabenhefte)
\newtcolorbox{loesung}[1][]{%
  enhanced, breakable,
  colback=bfwBgAmber!40,
  colframe=bfwSuccess,
  boxrule=0.6pt,
  arc=3pt,
  left=10pt, right=10pt, top=8pt, bottom=8pt,
  fonttitle=\bfseries\color{bfwSuccess!70!black},
  title={Lösung},
  #1
}

% Glossar-Eintrag
\newcommand{\glossareintrag}[2]{%
  \par\noindent\textbf{\color{bfwPrimaryDark}#1} \enspace #2\par\smallskip
}

% Section-Stil — etwas farbiger als Standard
\usepackage{titlesec}
\titleformat{\section}
  {\Large\bfseries\color{bfwPrimaryDark}}
  {\thesection}{0.6em}{}
\titleformat{\subsection}
  {\large\bfseries\color{bfwPrimary}}
  {\thesubsection}{0.5em}{}
\titleformat{\subsubsection}
  {\normalsize\bfseries\color{bfwInk}}
  {\thesubsubsection}{0.4em}{}

% TikZ-Stil "skizzenhaft" — etwas weicher als Rough.js, aber stabil
\tikzset{
  roughbox/.style = {
    draw=bfwPrimaryDark, thick, fill=bfwBgTeal,
    rounded corners=4pt, inner sep=6pt
  },
  roughaccent/.style = {
    draw=bfwAccent, thick, fill=bfwBgAmber,
    rounded corners=4pt, inner sep=6pt
  },
  roughline/.style = {
    draw=bfwInkSoft, thick, -{Stealth[length=2.5mm]}
  }
}

% Bullet-Symbole (bewusst kein FontAwesome — vermeidet Font-Installations-Probleme)
\newcommand{\faLightbulb}{$\bullet$}
\newcommand{\faBookmark}{$\bullet$}

% Header/Footer
\usepackage{fancyhdr}
\pagestyle{fancy}
\fancyhf{}
\renewcommand{\headrulewidth}{0pt}
\fancyfoot[L]{\small\color{bfwInkSoft}\thepage}
\fancyfoot[R]{\small\color{bfwInkSoft} BFW M-064 Beschaffung im Büromanagement}


\title{\vspace*{-1.5cm}\Huge\bfseries\color{bfwPrimaryDark}Tag~10: Postbearbeitung --- Posteingang\\[0.4em]
  \large\color{bfwInkSoft}\textnormal{Vom Briefkasten zum richtigen Schreibtisch}}
\author{\textsf{Büroprozesse gestalten und Arbeitsvorgänge organisieren \textperiodcentered{} M-067}\\
\textsf{\small\copyright\ Can Siebert\,\textperiodcentered\,2026}}
\date{10.07.2026}

\begin{document}
\maketitle
\thispagestyle{empty}

\vspace{1em}
\begin{merksatz}[title={\bfseries Worum geht es heute?}]
Jeder Brief, jede Rechnung und jeder Vertrag, der in ein Unternehmen gelangt, durchläuft
zuerst den \emph{Posteingang}. Diese Stelle ist der erste Umschlagplatz für Informationen
--- arbeitet sie sorgfältig, gelangt die richtige Post schnell zur richtigen Person; arbeitet
sie nachlässig, gehen Dokumente verloren, werden Fristen versäumt oder das Briefgeheimnis
verletzt. In diesem Handout lernen Sie, auf welchen Wegen Post ankommt, welche Arbeitsschritte
beim Posteingang in welcher Reihenfolge ablaufen, welche Post nicht geöffnet werden darf und
wie ein moderner, digitaler Posteingang funktioniert.
\end{merksatz}

\tableofcontents
\vspace{1em}
\hrule
\vspace{1em}

\section{Bedeutung des Posteingangs}

Die Stelle in einem Unternehmen, die die Post bearbeitet, ist der \emph{erste Umschlagplatz
für ein- und ausgehende Informationen} und spielt deshalb eine bedeutende Rolle. Über sie
laufen Aufträge, Rechnungen, Verträge und Behördenschreiben --- also alles, was geschäftlich
wichtig ist.

In \textbf{kleineren Unternehmen} wird die Bearbeitung der Tagespost meist vom
\emph{Sekretariat} erledigt: Die eingehende Post wird angenommen, sortiert, geöffnet und
verteilt; die ausgehende Post wird kuvertiert, frankiert und verschickt. In \textbf{größeren
Unternehmen} übernimmt diese Aufgaben eine eigene Abteilung, die \emph{Poststelle}.

\begin{merksatz}
Man unterscheidet die \textbf{zentrale Poststelle} (eine Stelle bearbeitet die gesamte Post
des Unternehmens --- einheitlich, spezialisiert, mit effizienter Technik) und die
\textbf{dezentrale Poststelle} (jede Abteilung bearbeitet ihre Post selbst --- kurze Wege,
hohe Sachnähe, aber mehrfache Ausstattung).
\end{merksatz}

\section{Möglichkeiten des Posteingangs}

Die Post kann auf verschiedene Arten das Unternehmen erreichen.

\subsection{Zustellung durch Briefzusteller / Postbote}

Laut Postgesetz ist die \textbf{Deutsche Post AG} dazu verpflichtet und berechtigt, Briefe,
Päckchen und Pakete dem Empfänger zuzustellen. Um unabhängig von den Zustellzeiten eines
Postboten zu sein, kann in einer Postfiliale ein \textbf{Postfach} gemietet werden; hierfür
berechnet die Post eine einmalige Einrichtungspauschale. Die Größe des Postfachs hängt vom
täglichen Postaufkommen des Unternehmens ab. Ein weiterer Service ist der \textbf{Service
Hin + Weg}: Die Deutsche Post AG bringt die Briefe zu vereinbarten Zeiten in das Unternehmen
und holt die ausgehenden Briefe auch dort ab.

\subsection{Zustellung durch einen privaten Dienstleister}

Mit der Aufhebung des Briefmonopols ist seit dem \textbf{1.\ Januar 2008} der deutsche
Briefmarkt gesetzlich liberalisiert. Nicht nur die Deutsche Post AG darf Briefe befördern,
sondern auch andere privatwirtschaftliche Unternehmen, die hierfür eine \emph{Lizenz der
Bundesnetzagentur} benötigen. Bekannte internationale private Brief- und Paketdienstleister
sind z.\,B.\ TNT, FedEx und UPS.

\subsection{Weitere Eingangswege}

Neben dem klassischen Briefweg erreicht ein Unternehmen Post heute auch über
\textbf{Fax}, \textbf{E-Mail}, die \textbf{persönliche Abgabe} am Empfang sowie durch
\textbf{Boten} (z.\,B.\ Kuriere oder Fahrer anderer Firmen). Auch diese Eingänge müssen erfasst
und an die richtige Stelle weitergeleitet werden.

\begin{eselsbruecke}
\textbf{„Post kommt per Bote, Brief, Fax und E-Mail.”} Merken Sie sich: Posteingang ist heute
\emph{mehr} als der Briefkasten --- der digitale Eingang (E-Mail, E-Rechnung) wird ständig
wichtiger.
\end{eselsbruecke}

\section{Arbeitsschritte beim Posteingang}

Um die Geschäftspost überhaupt entgegennehmen zu können, wird eine rechtsgültige
\textbf{Postvollmacht} benötigt. Dies ist eine einfache schriftliche Erklärung mit folgenden
Angaben: \emph{Vollmachtgeber, Bevollmächtigter, Gültigkeitsdauer, Umfang der Vollmacht}
(normale Post, Einschreiben, eigenhändig), \emph{Datum der Ausstellung} und \emph{Unterschrift
des Vollmachtgebers}.

Danach läuft die Bearbeitung in einer festen Reihenfolge ab:

\begin{enumerate}
  \item \textbf{Annehmen} --- die Sendungen werden vom Zusteller oder Dienstleister entgegengenommen.
  \item \textbf{Vorsortieren} --- die Sendungen werden \emph{ungeöffnet} vorsortiert.
  \item \textbf{Öffnen} --- die Geschäftspost wird geöffnet (private/vertrauliche Post bleibt zu).
  \item \textbf{Kontrollieren} --- Leerkontrolle und Anlagenkontrolle.
  \item \textbf{Stempeln} --- der Posteingangsstempel wird angebracht.
  \item \textbf{Verteilen} --- die Post gelangt zu den zuständigen Mitarbeitern.
\end{enumerate}

\subsection{Vorsortieren}

Zunächst werden die eingehenden Postsendungen \emph{ungeöffnet} vorsortiert in:
Geschäftsbriefe, private Briefe, Post an die Geschäftsleitung, an die Personalabteilung und
den Betriebsrat sowie \textbf{Irrläufer}. Irrläufer sind fehlgeleitete Sendungen, die
ungeöffnet dem Postzusteller oder dem privaten Dienstleister zurückgegeben werden.

\subsection{Öffnen der Post}

Private Briefe sowie Post an die Geschäftsleitung, die Personalabteilung und den Betriebsrat
werden \emph{ungeöffnet} weitergeleitet. \textbf{Private Post} erkennt man daran, dass im
Anschriftenfeld der \emph{Name einer natürlichen Person zuerst} genannt wird und erst dann die
Firma folgt; in der Zusatz- und Vermerkzone steht zusätzlich der Hinweis „Persönlich” oder
„Vertraulich”. Wurde private Post versehentlich geöffnet, sollte sie wieder verschlossen und
mit dem Vermerk „versehentlich geöffnet” weitergeleitet werden.

\begin{merksatz}
\textbf{Faustregel zum Öffnen:} Steht im Anschriftenfeld \emph{zuerst die Firma} und dann der
Name, handelt es sich um \emph{Geschäftspost} --- sie darf geöffnet und an die genannte Person
weitergeleitet werden. Steht zuerst der \emph{Name einer Privatperson} und es fehlt der Zusatz
„Persönlich/Vertraulich”, darf der Brief laut Urteil (LAG~Hamm 2003, AZ:~14 Sa 1972/02) in der
Poststelle ebenfalls geöffnet werden.
\end{merksatz}

\begin{center}
\begin{tabular}{p{2.6cm} p{2.9cm} p{2.7cm} p{2.7cm}}
\toprule
\textbf{Geschäftsbrief} & \textbf{Brief an die Geschäftsleitung} & \textbf{Privatbrief} & \textbf{Irrläufer}\\
\midrule
BüroKonzept KG\newline Herrn Achim Langner\newline Rochusstr.\ 30\newline 53123 Bonn
& BüroKonzept KG\newline Geschäftsleitung\newline Herrn Sebastian Falo\newline 53123 Bonn
& Persönlich\newline Frau Eva Holländer\newline BüroKonzept KG\newline 53123 Bonn
& Büro GmbH\newline Herrn Volker Berg\newline Rochusstr.\ 320\newline 53123 Bonn\\
\midrule
darf geöffnet werden & ungeöffnet weiterleiten & nicht öffnen & ungeöffnet zurück\\
\bottomrule
\end{tabular}
\end{center}

\medskip
Zum Öffnen der Geschäftspost wird entweder ein \emph{Handbrieföffner} (bei geringem
Postaufkommen) oder ein \emph{automatischer Brieföffner} verwendet.

\subsection{Kontrollieren}

Nachdem die Geschäftspost geöffnet und der Inhalt entnommen wurde, sollte kontrolliert werden,
ob der Umschlag auch vollständig geleert wurde (\textbf{Leerkontrolle}) und ob alle im Brief
genannten Anlagen beigefügt sind (\textbf{Anlagenkontrolle}). Fehlt eine Anlage, wird dies
handschriftlich auf dem Brief vermerkt.

Fehlt die Absenderangabe, handelt es sich um amtlich zugestellte Sendungen, oder liegen das
Datum des Briefes und der Poststempel zeitlich ungewöhnlich weit auseinander, so wird der
Briefumschlag als Nachweis an das Schriftstück angeheftet.

\subsection{Stempeln}

Ein \textbf{Posteingangsstempel} wird rechts oben neben die Anschrift des Empfängers angebracht.
Er enthält das aktuelle Datum, evtl.\ die Uhrzeit und Platz für Bearbeitungs- und
Erledigungsvermerke. \emph{Nicht} gestempelt werden Dokumente wie z.\,B.\ Urkunden, Zeugnisse,
Verträge, Schecks und Zeichnungen; hier empfiehlt es sich, den Umschlag zu stempeln und
beizufügen.

Der Posteingangsstempel ermöglicht einen Abgleich zwischen dem Datum des Schriftstücks und dem
Eingangsdatum, um Ursachen für eine verzögerte Postzustellung prüfen zu können. Außerdem lässt
sich mit seiner Hilfe die innerbetriebliche Bearbeitungszeit feststellen.

\subsection{Verteilen}

Nachdem die Eingangspost mit einem Eingangsstempel versehen wurde, kann sie an die
entsprechenden Mitarbeiter in den jeweiligen Abteilungen verteilt werden. Dafür gibt es
mehrere Möglichkeiten:

\begin{itemize}
  \item \textbf{Ablage in Mitarbeiter- oder Abteilungsfächer}, aus denen die Post abgeholt wird.
  \item \textbf{Ein Bote} übernimmt die Verteilung in die verschiedenen Abteilungen.
  \item \textbf{Automatische Verteilung} über Schriftgutförderanlagen, z.\,B.\ ein Rohrpostsystem.
\end{itemize}

\subsection{Behandlung von Wertsendungen}

\textbf{Wertsendungen} wie Einschreiben, Wertbriefe oder Nachnahmesendungen erfordern besondere
Sorgfalt: Sie werden gegen Quittung angenommen (für Einschreiben und eigenhändige Sendungen ist
die Postvollmacht nötig), in einem \emph{Eingangs- bzw.\ Postbuch} erfasst und sicher verwahrt
oder unverzüglich gegen Nachweis an den benannten Empfänger übergeben. So bleibt jederzeit
nachvollziehbar, wer eine wichtige Sendung erhalten hat.

\section{Posteingangssysteme und digitale Archivierung}

Unternehmen mit hohem Postaufkommen nutzen \textbf{Posteingangssysteme}, bei denen die
verschiedenen Arbeitsschritte bei der Bearbeitung des Posteingangs automatisiert erfolgen:

\begin{enumerate}
  \item maschinelle Öffnung der Briefumschläge mithilfe von elektrischen Brieföffnern;
  \item automatische Entnahme des Inhalts und elektronische Leerkontrolle des Briefumschlags;
  \item Einlesen der Schriftstücke mithilfe von \emph{Scannern} und Speicherung im
    unternehmensinternen EDV-System. Die Mitarbeiter können über ihren eigenen PC direkt auf
    die für sie bestimmte Post zugreifen.
\end{enumerate}

In der modernen Praxis wird beim Scannen häufig eine \textbf{OCR-Texterkennung} (Optical
Character Recognition) eingesetzt: Der eingescannte Brief wird in durchsuchbaren Text
umgewandelt und in einem Dokumentenmanagementsystem (DMS) abgelegt, oft mit automatischer
Zuordnung zum richtigen Vorgang. Damit verschmelzen Papier- und digitaler Posteingang.

\begin{merksatz}
Die Digitalisierung des Posteingangs hat für ein Unternehmen viele Vorteile: schnellerer
Informationsfluss; keine Zeitverzögerung durch aufwändige Verteilung (kürzere Bearbeitungszeit);
Schriftstücke sind jederzeit und schnell auffindbar; der Papierverbrauch sinkt; die
Auskunftsfähigkeit steigt; die Post kann auf ihrem Weg nicht mehr verloren gehen; und die
Effizienz der internen Prozessabläufe steigt.
\end{merksatz}

\section{Zusammenfassung}

\begin{enumerate}
  \item Der Posteingang ist der \textbf{erste Umschlagplatz für Informationen}; in kleinen Firmen erledigt ihn das Sekretariat, in großen die Poststelle (zentral oder dezentral).
  \item Post erreicht das Unternehmen über \textbf{Briefzusteller, Postfach, Service Hin + Weg, private Dienstleister} (seit 2008 liberalisiert) sowie \textbf{Fax, E-Mail, Boten und persönliche Abgabe}.
  \item Die \textbf{Arbeitsschritte} lauten: Annehmen $\rightarrow$ Vorsortieren $\rightarrow$ Öffnen $\rightarrow$ Kontrollieren $\rightarrow$ Stempeln $\rightarrow$ Verteilen; Voraussetzung ist die \textbf{Postvollmacht}.
  \item \textbf{Privat-/vertrauliche Post} wird nicht geöffnet; steht die Firma zuerst, ist es Geschäftspost und darf geöffnet werden.
  \item Der \textbf{Posteingangsstempel} dokumentiert Eingangsdatum und Bearbeitungszeit; \textbf{Wertsendungen} werden gesondert erfasst und verwahrt.
  \item \textbf{Posteingangssysteme} mit \textbf{Scannen/OCR} und digitaler Archivierung beschleunigen den Informationsfluss und machen Schriftstücke jederzeit auffindbar.
\end{enumerate}

\newpage

\section*{Glossar}
\addcontentsline{toc}{section}{Glossar}

\glossareintrag{Posteingang}{Erster Umschlagplatz für eingehende Informationen; Annahme und Bearbeitung der eingehenden Post.}

\glossareintrag{Poststelle}{Eigene Abteilung größerer Unternehmen, die die gesamte Postbearbeitung übernimmt.}

\glossareintrag{Zentrale Poststelle}{Eine Stelle bearbeitet die gesamte Post des Unternehmens; einheitlich und spezialisiert.}

\glossareintrag{Dezentrale Poststelle}{Jede Abteilung bearbeitet ihre Post selbst; kurze Wege, hohe Sachnähe.}

\glossareintrag{Postfach}{In einer Postfiliale gemieteter Behälter; macht unabhängig von den Zustellzeiten des Postboten.}

\glossareintrag{Service Hin + Weg}{Service der Deutschen Post AG: Briefe werden zu vereinbarten Zeiten gebracht und abgeholt.}

\glossareintrag{Bundesnetzagentur}{Behörde, die privaten Brief- und Paketdienstleistern die nötige Lizenz erteilt.}

\glossareintrag{Postvollmacht}{Rechtsgültige schriftliche Erklärung, um Geschäftspost (auch Einschreiben/eigenhändig) entgegennehmen zu können.}

\glossareintrag{Irrläufer}{Fehlgeleitete Sendung, die ungeöffnet dem Zusteller oder Dienstleister zurückgegeben wird.}

\glossareintrag{Leerkontrolle}{Prüfung, ob der Briefumschlag nach Entnahme des Inhalts vollständig geleert wurde.}

\glossareintrag{Anlagenkontrolle}{Prüfung, ob alle im Brief genannten Anlagen beigefügt sind; fehlende werden vermerkt.}

\glossareintrag{Posteingangsstempel}{Stempel rechts oben neben der Anschrift mit Datum, Uhrzeit und Bearbeitungsvermerken.}

\glossareintrag{Wertsendung}{Einschreiben, Wertbrief oder Nachnahme; wird gesondert erfasst, quittiert und sicher verwahrt.}

\glossareintrag{Schriftgutförderanlage}{Technik zur automatischen Verteilung der Post, z.\,B.\ ein Rohrpostsystem.}

\glossareintrag{Posteingangssystem}{System, das die Arbeitsschritte des Posteingangs automatisiert (Öffnen, Entnahme, Scannen).}

\glossareintrag{OCR}{Optical Character Recognition; Texterkennung, die gescannte Dokumente durchsuchbar macht.}

\glossareintrag{Digitale Archivierung}{Speicherung der eingescannten Eingangspost im EDV-System; jederzeit auffindbar.}

\end{document}
